% Figures holding a centred TikZ picture of a fixed size, as in real papers: in center, 0.525 pt wider than the text (TeX reports the line overfull and sets the picture on a line of its own, \lineskip below an empty first line), one narrower, and one wider with \centering, each with a caption and text after it.
\documentclass[11pt]{article}
\usepackage{tikz}
% one page with nothing but the text: checked against plain lualatex too
% (cases.ts), where every float is [H]
\pagestyle{empty}
\newcommand\testpicture[1]{%
  \begin{tikzpicture}
    \useasboundingbox (0,0) rectangle (#1,2cm);
    \draw (0.2,0.2) rectangle (2,1.8);
    \node at (1.1,1) {$a_1$};
    \node[anchor=east] at (#1,1) {right};
  \end{tikzpicture}}
\begin{document}
A paragraph before the first figure, long enough to break onto a second line
of the column, quickly.
\begin{figure}
\begin{center}
  \testpicture{360.525pt}
\end{center}
\caption{A picture a little wider than the text, in center.}
\end{figure}
A paragraph after the first figure, long enough to break onto a second line
of the column, before the next one.
\begin{figure}
\begin{center}
  \testpicture{330pt}
\end{center}
\caption{A picture narrower than the text, in center.}
\end{figure}
A paragraph after the second figure, long enough to break onto a second line
of the column, before the last one.
\begin{figure}
  \centering
  \testpicture{360.525pt}
  \caption{A picture a little wider than the text, centred with centering.}
\end{figure}
A last paragraph, after all the figures, long enough to break onto a second
line of the column.
\end{document}
